% GENERATED FILE. Edit canonical lessons, then run npm run generate:books.
% canonical-source-hash: fnv1a64:b2cc679d2c60d2e7
% canonical-lessons: PA-C221-future-return, PA-C221-javanga, PA-C221-javanga-review, PA-C221-javega, PA-C221-javange, PA-C221-payoff

\chapter{I Will Go}
\label{ch:pa-a2-javanga}
\hlchaptermodality{221}

\begin{chapteropening}
\textbf{By the end of this chapter:} \emph{I can say that I, someone else or we will go somewhere tomorrow, with the -ga future.}

Say who will go where tomorrow with javanga / javangi, javega / javegi and javange.
\end{chapteropening}

\section[\texorpdfstring{\pa{ਹੋਵੇਗਾ} · \pa{ਆਵੇਗਾ} … (hovegā · āvegā …)}{ਹੋਵੇਗਾ · ਆਵੇਗਾ … (hovegā · āvegā …)}]{Return to three familiar futures}

\label{lesson:PA-C221-future-return}



\begin{quote}
\emph{Hear:} \emph{hovegā; āvegā} Say \textbf{\pa{ਹੋਵੇਗਾ} · \pa{ਆਵੇਗਾ}}, \textquotedblleft{}will be; will come\textquotedblright{}. Then the farewell \textbf{\pa{ਫਿਰ} \pa{ਮਿਲਾਂਗੇ}}, \textquotedblleft{}we will meet again\textquotedblright{}.
\end{quote}

\begin{grammarlens}[title={Grammar Lens}]
All three end in \textbf{-\pa{ਗਾ}} or \textbf{-\pa{ਗੇ}}, the future. \textbf{\pa{ਹੋਵੇਗੀ}} and \textbf{\pa{ਆਵੇਗੀ}} have \textbf{-\pa{ਗੀ}} for a feminine subject. \textbf{\pa{ਮਿਲਾਂਗੇ}} is \textquotedblleft{}we will meet\textquotedblright{}. The next lesson gives you \emph{I will go}.
\end{grammarlens}

\subsection*{Your turn}
\begin{itemize}
\raggedright
  \item \emph{Hear:} \emph{Aman will come} Say \textbf{\pa{ਅਮਨ} \pa{ਆਵੇਗਾ}}.
  \item \emph{Hear:} \emph{the train will come} Say \textbf{\pa{ਰੇਲ} \pa{ਆਵੇਗੀ}}.
  \item \emph{Hear:} \emph{we will meet tomorrow} Say \textbf{\pa{ਕੱਲ੍ਹ} \pa{ਮਿਲਾਂਗੇ}}.
\end{itemize}

\begin{tcolorbox}[breakable,colback=teal!4,colframe=teal!35!black,title={Before you move on}]
Say \emph{will come}, then \emph{we will meet}.
\end{tcolorbox}
\section[\texorpdfstring{\pa{ਜਾਵਾਂਗਾ} · \pa{ਜਾਵਾਂਗੀ} (jāvāṅgā · jāvāṅgī)}{ਜਾਵਾਂਗਾ · ਜਾਵਾਂਗੀ (jāvāṅgā · jāvāṅgī)}]{\pa{ਜਾਵਾਂਗਾ}, \pa{ਜਾਵਾਂਗੀ} — I will go}

\label{lesson:PA-C221-javanga}



\begin{quote}
Say \emph{Aman will come}. Then \emph{we will meet tomorrow}.
\end{quote}

\begin{grammarlens}[title={Grammar Lens: \pa{ਮੈਂ} in the future}]
After \textbf{\pa{ਜਾ}} comes \textbf{\pa{ਵ}}, as in \textbf{\pa{ਆਵੇਗਾ}}. With \textbf{\pa{ਮੈਂ}} the ending is \textbf{-\pa{ਾਂਗਾ}} for a man and \textbf{-\pa{ਾਂਗੀ}} for a woman:

\begin{quote}
\textbf{\pa{ਜਾ}} + \textbf{\pa{ਵ}} + \textbf{\pa{ਾਂਗਾ}} $\to$ \textbf{\pa{ਜਾਵਾਂਗਾ}}
\end{quote}

\begin{quote}
\textbf{\pa{ਮੈਂ} \pa{ਕੱਲ੍ਹ} \pa{ਸਕੂਲ} \pa{ਜਾਵਾਂਗਾ।}} — \emph{maiṁ kallh sakūl jāvāṅgā.} — I will go to school tomorrow (a man speaking).
\end{quote}

\begin{quote}
\textbf{\pa{ਮੈਂ} \pa{ਕੱਲ੍ਹ} \pa{ਸਕੂਲ} \pa{ਜਾਵਾਂਗੀ।}} — \emph{maiṁ kallh sakūl jāvāṅgī.} — I will go to school tomorrow (a woman speaking).
\end{quote}

With a future form, \textbf{\pa{ਕੱਲ੍ਹ}} means tomorrow.
\end{grammarlens}

\subsection*{Your turn}
\begin{itemize}
\raggedright
  \item \emph{Hear:} \emph{I will go home tomorrow} Say it with the ending that fits you.
  \item \emph{Hear:} \emph{I will go to the market} Say \textbf{\pa{ਮੈਂ} \pa{ਬਾਜ਼ਾਰ} \pa{ਜਾਵਾਂਗਾ।}} or \textbf{\pa{ਮੈਂ} \pa{ਬਾਜ਼ਾਰ} \pa{ਜਾਵਾਂਗੀ।}}
  \item \emph{Read it:} \textbf{\pa{ਜਾਵਾਂਗਾ}} Point to \textbf{\pa{ਵ}}, then to \textbf{\pa{ਗਾ}}.
  \item \emph{Write it:} copy \textbf{\pa{ਜਾਵਾਂਗਾ}} once
\end{itemize}

\begin{tcolorbox}[breakable,colback=teal!4,colframe=teal!35!black,title={Before you move on}]
\emph{I will go}, as a man? (\textbf{\pa{ਜਾਵਾਂਗਾ}}.) As a woman? (\textbf{\pa{ਜਾਵਾਂਗੀ}}.)
\end{tcolorbox}
\section[\texorpdfstring{\pa{ਜਾਵਾਂਗਾ} · \pa{ਜਾਵਾਂਗੀ} (jāvāṅgā · jāvāṅgī)}{ਜਾਵਾਂਗਾ · ਜਾਵਾਂਗੀ (jāvāṅgā · jāvāṅgī)}]{Hear and say I will go again}

\label{lesson:PA-C221-javanga-review}



\begin{quote}
\emph{Hear:} \emph{jāvāṅgā; jāvāṅgī} Say \textbf{\pa{ਜਾਵਾਂਗਾ} · \pa{ਜਾਵਾਂਗੀ}}.
\end{quote}

\subsection*{Your turn}
\begin{itemize}
\raggedright
  \item \emph{Hear:} \emph{I will go to school tomorrow} Say it with the ending that fits you.
  \item \emph{Hear:} \emph{I went to school yesterday} Say \textbf{\pa{ਮੈਂ} \pa{ਕੱਲ੍ਹ} \pa{ਸਕੂਲ} \pa{ਗਿਆ।}} or \textbf{\pa{ਮੈਂ} \pa{ਕੱਲ੍ਹ} \pa{ਸਕੂਲ} \pa{ਗਈ।}}
  \item \emph{Read it:} \textbf{\pa{ਮੈਂ} \pa{ਕੱਲ੍ਹ} \pa{ਘਰ} \pa{ਜਾਵਾਂਗੀ।}} Give its meaning.
  \item \emph{Write it:} cover the model and write \textbf{\pa{ਜਾਵਾਂਗੀ}} once
\end{itemize}

\begin{tcolorbox}[breakable,colback=teal!4,colframe=teal!35!black,title={Before you move on}]
Same \textbf{\pa{ਕੱਲ੍ਹ}}, two days: which verb makes it tomorrow? (The future.)
\end{tcolorbox}
\section[\texorpdfstring{\pa{ਜਾਵੇਗਾ} · \pa{ਜਾਵੇਗੀ} (jāvegā · jāvegī)}{ਜਾਵੇਗਾ · ਜਾਵੇਗੀ (jāvegā · jāvegī)}]{\pa{ਜਾਵੇਗਾ}, \pa{ਜਾਵੇਗੀ} — will go}

\label{lesson:PA-C221-javega}



\begin{quote}
Say \emph{I will go tomorrow}. Then \emph{Aman will come} and \emph{the train will come}.
\end{quote}

\begin{grammarlens}[title={Grammar Lens: like \pa{ਆਵੇਗਾ}}]
For one other person the ending is the one \textbf{\pa{ਆਵੇਗਾ}} has: \textbf{-\pa{ੇਗਾ}} for a man, \textbf{-\pa{ੇਗੀ}} for a woman.

\begin{quote}
\textbf{\pa{ਅਮਨ} \pa{ਕੱਲ੍ਹ} \pa{ਘਰ} \pa{ਜਾਵੇਗਾ।}} — \emph{Aman kallh ghar jāvegā.} — Aman will go home tomorrow.
\end{quote}

\begin{quote}
\textbf{\pa{ਕੁੜੀ} \pa{ਕੱਲ੍ਹ} \pa{ਸਕੂਲ} \pa{ਜਾਵੇਗੀ।}} — \emph{kuṛī kallh sakūl jāvegī.} — The girl will go to school tomorrow.
\end{quote}
\end{grammarlens}

\subsection*{Your turn}
\begin{itemize}
\raggedright
  \item \emph{Hear:} \emph{he will go to the market} Say \textbf{\pa{ਉਹ} \pa{ਬਾਜ਼ਾਰ} \pa{ਜਾਵੇਗਾ।}}
  \item \emph{Hear:} \emph{the woman will go home tomorrow} Say \textbf{\pa{ਔਰਤ} \pa{ਕੱਲ੍ਹ} \pa{ਘਰ} \pa{ਜਾਵੇਗੀ।}}
  \item \emph{Read it:} \textbf{\pa{ਅਮਨ} \pa{ਕੱਲ੍ਹ} \pa{ਘਰ} \pa{ਜਾਵੇਗਾ।}} Give its meaning.
  \item \emph{Write it:} copy \textbf{\pa{ਜਾਵੇਗਾ}} once
\end{itemize}

\begin{tcolorbox}[breakable,colback=teal!4,colframe=teal!35!black,title={Before you move on}]
\emph{He will go}? (\textbf{\pa{ਜਾਵੇਗਾ}}.) \emph{She will go}? (\textbf{\pa{ਜਾਵੇਗੀ}}.)
\end{tcolorbox}
\section[\texorpdfstring{\pa{ਅਸੀਂ} \pa{ਜਾਵਾਂਗੇ} (asīṁ jāvāṅge)}{ਅਸੀਂ ਜਾਵਾਂਗੇ (asīṁ jāvāṅge)}]{\pa{ਅਸੀਂ} \pa{ਜਾਵਾਂਗੇ} — we will go}

\label{lesson:PA-C221-javange}



\begin{quote}
Say \emph{Aman will go home tomorrow}. Then \emph{we will meet again}.
\end{quote}

\begin{grammarlens}[title={Grammar Lens: the ending of \pa{ਮਿਲਾਂਗੇ}}]
\textbf{\pa{ਮਿਲਾਂਗੇ}} has been \textquotedblleft{}we will meet\textquotedblright{} since the early lessons. Its ending, \textbf{-\pa{ਾਂਗੇ}}, is the one for \textbf{\pa{ਅਸੀਂ}}:

\begin{quote}
\textbf{\pa{ਅਸੀਂ} \pa{ਕੱਲ੍ਹ} \pa{ਬਾਜ਼ਾਰ} \pa{ਜਾਵਾਂਗੇ।}} — \emph{asīṁ kallh bāzār jāvāṅge.} — We will go to the market tomorrow.
\end{quote}

\begin{quote}
\textbf{\pa{ਅਸੀਂ} \pa{ਸਵੇਰੇ} \pa{ਜਾਵਾਂਗੇ।}} — \emph{asīṁ savere jāvāṅge.} — We will go in the morning.
\end{quote}
\end{grammarlens}

\subsection*{Your turn}
\begin{itemize}
\raggedright
  \item \emph{Hear:} \emph{we will go home tomorrow} Say \textbf{\pa{ਅਸੀਂ} \pa{ਕੱਲ੍ਹ} \pa{ਘਰ} \pa{ਜਾਵਾਂਗੇ।}}
  \item \emph{Hear:} \emph{I will go, and Aman will go} Say \textbf{\pa{ਮੈਂ} \pa{ਜਾਵਾਂਗਾ} \pa{ਅਤੇ} \pa{ਅਮਨ} \pa{ਜਾਵੇਗਾ।}}
  \item \emph{Read it:} \textbf{\pa{ਅਸੀਂ} \pa{ਸਵੇਰੇ} \pa{ਜਾਵਾਂਗੇ।}} Give its meaning.
  \item \emph{Write it:} copy \textbf{\pa{ਜਾਵਾਂਗੇ}} once
\end{itemize}

\begin{tcolorbox}[breakable,colback=teal!4,colframe=teal!35!black,title={Before you move on}]
\emph{We will go}? (\textbf{\pa{ਅਸੀਂ} \pa{ਜਾਵਾਂਗੇ}}.) \emph{We will meet}? (\textbf{\pa{ਅਸੀਂ} \pa{ਮਿਲਾਂਗੇ}}.)
\end{tcolorbox}
\section[\texorpdfstring{\pa{ਮੈਂ} \pa{ਕੱਲ੍ਹ} \pa{ਸਕੂਲ} \pa{ਜਾਵਾਂਗਾ।} (maiṁ kallh sakūl …)}{ਮੈਂ ਕੱਲ੍ਹ ਸਕੂਲ ਜਾਵਾਂਗਾ। (maiṁ kallh sakūl …)}]{Tomorrow: a four-skill payoff}

\label{lesson:PA-C221-payoff}



\begin{quote}
Say \emph{I will go}, \emph{she will go}, \emph{we will go}.
\end{quote}

\subsection*{Your turn: four skills}
\begin{itemize}
\raggedright
  \item \emph{Hear:} \emph{I will go to the market tomorrow} Say it with the ending that fits you.
  \item \emph{Hear:} \emph{she went yesterday; he will go tomorrow} Say \textbf{\pa{ਉਹ} \pa{ਕੱਲ੍ਹ} \pa{ਗਈ।} \pa{ਉਹ} \pa{ਕੱਲ੍ਹ} \pa{ਜਾਵੇਗਾ।}}
  \item \emph{Read it:} \textbf{\pa{ਅਸੀਂ} \pa{ਕੱਲ੍ਹ} \pa{ਸਕੂਲ} \pa{ਜਾਵਾਂਗੇ।}} Give its meaning: we will go to school tomorrow.
  \item \emph{Write it:} cover the line and write \textbf{\pa{ਅਮਨ} \pa{ਕੱਲ੍ਹ} \pa{ਘਰ} \pa{ਜਾਵੇਗਾ।}}
\end{itemize}

\begin{tcolorbox}[breakable,colback=teal!4,colframe=teal!35!black,title={Before you move on}]
Say where you will go tomorrow, and who will go with you.
\end{tcolorbox}
